% Modernised reading — fonds Alexandre Grothendieck, shelfmark 129
% (pages 1-114, the whole folder).
%
% This is an INTERPRETATION, not a transcription. It re-reads the pages in
% today's notation and vocabulary, states the hypotheses the manuscript leaves
% implicit, and drops the critical apparatus so that the mathematics reads
% straight through. Everything the brackets and underlines carried moves into a
% footnote. It is held to being mathematically correct as it stands; where the
% manuscript is loose or mistaken, the true statement is what appears here and
% the footnote says what the page has. In any dispute of substance the
% transcription governs: whoever cites Grothendieck must cite batch-01.fr.tex
% (pp. 1-20) ... batch-06.fr.tex (pp. 101-114).
%
% Pass: Opus 5.5 (claude-opus-5-5), 2026-10-03 — first pass, unchecked against the pages by a human.
% The folder was read whole, its six batches together; this is one document
% for the shelfmark. It was made on Opus 5.5 and has not been measured against
% the reference reading of folder 115 (made on Opus 5). The literature named
% in the footnotes is cited from memory and was not checked against the
% sources; in particular the identification of pp. 112-114 with the 1970 Nice
% congress address is ours and was not compared with the printed text.
%
% Scope. Pages 1-75 are seminar typescripts by other people (chapter I
% attributed by a pencilled « DELALE » on p. 1; chapters II and III
% unsigned), with proof corrections in a hand nobody has identified. Per
% RIGHTS.md and the transcriptions' rule, they are identified here no further
% than the transcription's notes do; their mathematics is not restated. Pages
% 76 and 82 are listing-paper labels; 77, 79-80, 83-87, 89-90 are
% Grothendieck's own ribbon typescript (authorship argued in batch-05's
% header, internal evidence only, flagged there FOR THE HUMAN CHECK); 81,
% 91-95, 97, 99, 101-109, 112-114 are in his hand (106's first four English
% lines in an undetermined hand). Blank: 14, 78, 88, 111. Pages 98, 100, 102,
% 104, 108, 110 are duplicator copies of an unrelated text (« n° 380 / 382 »),
% not described. Pencil numbering jumps from 80 to 83 between PDF pages 80
% and 81 (batch-05 header): two leaves may be missing, and p. 80's sentence
% has no continuation in the folder.
%
% Spine. The Montreal course on Barsotti-Tate groups and Dieudonné crystals.
%   pp. 1-75    the seminar write-up, chapters I (Witt), II (finite groups,
%               classical Dieudonné), III (BT groups); identified only.
%   pp. 76-77   label; his typed plan in six chapters, renumbered in pencil.
%   pp. 79-81   chapters I-II in his draft: Witt vectors, W_->, K/W; the
%               decomposition over a perfect field and Dieudonné's theorem.
%   pp. 82-87   chapter III in his draft: examples, the connected-étale
%               filtration and its failure over a trait, Prop. (i)-(iv) with
%               its Lemma, extensions and cokernels, height.
%   pp. 89-90   groups with nilpotent Frobenius, Inf^infty(G) = G-hat (Thm
%               2.5); deformations of BT groups (Thm 4.1).
%   pp. 91-95   « Complexe cotangent relatif » (A)-(D): L, the co-Lie complex,
%               the « formule remarquable », link with Dieudonné, deformations
%               of flat groups, truncated BT groups, Ext over Lambda_n.
%   pp. 96-107  « Sorites »: Tate objects in an abelian category, then in
%               commutative group schemes; a subgroup of a level-2 object.
%   p. 109      formally smooth epimorphisms and a functor Phi.
%   pp. 112-114 « Conférence à Nice »: BT groups, Serre-Tate, PD, crystalline
%               site, crystals, programme I)-V).
%
% Notation fixed for the document.
%  - p a fixed prime; BT = Barsotti-Tate = p-divisible. G(n) = Ker(p^n . id_G)
%    (today G[p^n]), as in the pages; G[i] = Ker F^i (p. 89), never Ker p^i.
%  - D^* contravariant Dieudonné functor; F_M sigma-linear, V_M
%    sigma^{-1}-linear, FV = VF = p; Dieudonné ring D = W_sigma[F, V].
%  - ell^G co-Lie complex (homological degrees 0, 1), omega_G = H_0,
%    n_G = H_1, t_G = Lie(G), nu_G; the page's L^{G(n)} on p. 94 read ell.
%  - « objet de Tate de degré n »: (A, f), f^n = 0, Ker f^j = Im f^{n-j}; the
%    page writes « hauteur » and « degré » for the same n.
%  - On p. 112 « d » is the height (rank of G(1) = p^d), not a dimension.
%
% Substantive departures, each footnoted where it occurs:
%  p. 79   (x_i) -> (x_i^p) is a ring map only in char p (ours).
%  p. 80   W_-> is an ind-scheme; the bottom label read « r » taken as p
%          (T o F = p in char p, ours).
%  p. 81   second relation read VF = p (page: FV twice).
%  p. 83   « dont » read « donc »; the Ga-hat criterion left as his.
%  p. 84   « (et réciproquement) » read with his marginal.
%  pp. 85-86 the separable rank is that of G(1); the Lemma's factorisation is
%          radicial surjective THEN étale (the page types the reverse).
%  p. 89   target G[n-i]^{(p^i)} (page: exponent p^{n-i}); f read as the
%          relative Frobenius (ours).
%  p. 90   p locally nilpotent on S supplied for a), b); torsor under
%          t_{G0*} (x) t_{G0} (x) I (page omits I).
%  p. 92   the formula stated for G finite locally free; the marginal's
%          t_{G*} corrected to t_G to agree with the box.
%  p. 94   Ext of abelian sheaves (page indexes O_{S0}).
%  p. 97   rank supplied non-negative; (iv) only needs the inclusion (ours).
%  p. 101  A, B of the same degree n supplied; p. 101's head continues p. 99
%          (ours: p. 100 is a stray front).
%  pp. 105-107 generic fibre of G'_2 of height 2 supplied: without it a) does
%          not imply c) (Z/p inside Z/p^2, ours); alpha'' oriented G''_1 ->
%          G''_2 (the kernel and cokernel the page computes are of that map;
%          the page draws the verticals downward); setting (a trait) ours.
%  p. 112  Serre-Tate: p locally nilpotent, nilpotent ideal supplied.
%  p. 113  connection quasi-nilpotent supplied.
%  p. 114  omega_G read under the overwritten exponent; the fibre condition
%          is unnecessary in Messing's theorem.
%
% Links the pages do not write (ours, said so in the body): the seminar
% chapters I-III follow his typed plan and draft (by the notes' titles);
% pp. 97-107 with f = p are truncated BT groups, and p. 89 is their Frobenius
% analogue; p. 93's « Delta^*(G) » and p. 114's filtration are the Hodge
% filtration of the Dieudonné crystal; pp. 112-114 are an outline of the Nice
% address of 1970.
%
% Announced and never established: the continuation of p. 80; Notes I (props
% 1 and 14 cited p. 89); the conjectures of p. 83 (finite flatness of G(1),
% locally constant height); the proofs of Thm 2.5 and 4.1; the lettre à
% Illusie (cited, not in the folder); p. 93 entirely (conditional); (ii) and
% (ii bis) of 2.2; the proof of the Prop. of p. 94 and Th. of p. 95; the « 1st
% conj. » of p. 106; the dual list a*)-f*) of p. 107; most of p. 109 (Phi
% undefined); II)-V) of the Nice programme (« il se pourrait »).
\documentclass[11pt,a4paper]{article}
\input{../preamble/grothendieck}

\folder{129}
\pages{1}{114}
\dating{[1970]}
\watermark{Édition de démonstration\\interprétation personnelle de l'œuvre}
\foldertitle{Notes séminaire Montréal (Delale, Labute, Hakim, Messing) : copies de tapuscrits annotés (s.d.), notes manuscrites (s.d.) — lecture modernisée du dossier entier}

\begin{document}

\section*{Résumé}

\begin{resume}
Ce dossier est l'atelier d'un cours. Il rassemble, sous un titre de
l'inventaire qui nomme quatre auditeurs (Delale, Labute, Hakim, Messing),
d'une part trois chapitres de notes de séminaire mises au propre par
d'autres, d'autre part les brouillons de Grothendieck lui-même — un plan
tapé à la machine, des pages dactylographiées couvertes de corrections, des
feuilles manuscrites — et, pour finir, le canevas d'une conférence intitulée
« Conférence à Nice ». Le sujet est un seul objet : ce qu'on appelle
aujourd'hui un groupe $p$-divisible, et que ces pages appellent groupe de
Barsotti–Tate.

Le point de départ est familier. Sur un cercle, les racines de l'unité
d'ordre $p$, $p^{2}$, $p^{3}$… forment une tour de groupes finis emboîtés,
chacun obtenu du précédent en « divisant par $p$ ». Une courbe elliptique,
ou plus généralement une variété abélienne, porte une tour du même genre :
ses points de $p^{n}$-torsion. Cette tour retient une part étonnante de
l'information sur la variété, surtout lorsqu'on travaille en
caractéristique $p$, là où les points de torsion cessent d'être des points
et deviennent des objets « infinitésimaux ». Un groupe de Barsotti–Tate est
une telle tour, prise pour elle-même.

Le problème est de décrire ces objets par de l'algèbre linéaire. Sur un
corps parfait, Dieudonné l'avait fait : à chaque groupe correspond un
module sur un anneau de vecteurs de Witt, muni de deux opérateurs $F$ et
$V$ dont le produit est la multiplication par $p$. Les premières pages
rappellent cette théorie. L'idée nouvelle est de la faire vivre au-dessus
d'une base quelconque, en remplaçant le module par un \emph{cristal} : un
objet qui ne dépend pas d'un choix de relèvement de la base en
caractéristique nulle, mais qui sait se transporter rigidement le long de
tous ces relèvements à la fois. L'image est celle d'un cristal au sens
propre, dont la forme se propage identique à travers la matière. Grâce à
cette rigidité, déformer un groupe revient à déformer une donnée linéaire,
une filtration dans un module : le passage d'un problème non linéaire à un
problème d'algèbre linéaire est tout le gain.

Entre les deux, le dossier construit les outils. Il étudie comment un
groupe se déforme quand on épaissit la base, à l'aide du complexe
cotangent, qui mesure les obstructions comme le fait la dérivée pour une
fonction ; et il dégage, sous le nom d'« objets de Tate », la structure
abstraite d'un objet muni d'un endomorphisme nilpotent qui ressemble à un
module libre sur $\mathbb{Z}/p^{n}$. Le dossier ne démontre pas les
grands énoncés : il les ordonne, en signale les trous, et dit plusieurs
fois qu'il ne sait pas — « je ne sais pas procéder directement ! », « il se
pourrait que le foncteur soit pleinement fidèle, et j'ai une idée de ce que
pourrait être l'image essentielle ». C'est ce qui le rend lisible : un
programme écrit au moment où il se forme, et dont les étapes portent
aujourd'hui des noms.

La datation vient de l'inventaire ; aucune page n'est datée. Les noms sous
lesquels chercher la suite : groupes $p$-divisibles, théorie de Dieudonné
classique puis cristalline, cristaux et site cristallin, théorème de
Serre–Tate, théorème de Grothendieck–Messing sur les déformations,
complexe cotangent d'Illusie, filtration de Hodge du cristal de Dieudonné.

\keywords{p-divisible group, Barsotti-Tate group, truncated Barsotti-Tate group, Witt vectors, Dieudonné module, Dieudonné crystal, crystalline site, divided powers, Serre-Tate theorem, Grothendieck-Messing deformation theory, Hodge filtration, cotangent complex, co-Lie complex, deformation of group schemes, connected-étale sequence, formal Lie group, Cartier duality}
\end{resume}

\pagerange{1}{114}
\section*{Le fil du dossier, et les conventions}

\subsection*{Les stations}

Le dossier a deux strates, qu'on garde distinctes. Les pages 1 à 75 sont
des tapuscrits de séminaire dus à d'autres que lui ; les pages 76 à 114
sont de lui, en partie tapées à sa machine, en partie à la main. La lecture
ne porte que sur la seconde strate ; la première est seulement
identifiée\footnote{Les tapuscrits ne sont pas transcrits (voir l'en-tête
des lots 1 à 4), et le cadre juridique du projet interdit de restituer un
texte d'autrui au-delà de ce que disent les notes de la transcription. On
s'y tient.}.

\begin{enumerate}
  \item \emph{Pages 1 à 75.} Trois chapitres d'un tapuscrit de séminaire :
  « Préliminaires sur Witt » (deux frappes), « Groupes finis localement
  libres et théorie classique de Dieudonné », « Groupes de Barsotti–Tate ».
  \item \emph{Pages 76 et 77.} Une étiquette, puis le plan du cours en six
  chapitres, tapé et renuméroté au crayon.
  \item \emph{Pages 79 à 81.} Chapitres I et II du cours dans sa rédaction
  : vecteurs de Witt, l'ind-schéma $\underline{W}_{\to}$ et
  $K/W$ ; la décomposition d'un $p$-groupe fini sur un corps parfait et le
  théorème de Dieudonné.
  \item \emph{Pages 82 à 87.} Chapitre III, « Propriétés générales des
  B.T. » : exemples, filtration connexe-étale et son échec sur un trait,
  une proposition et son lemme, extensions et conoyaux, hauteur.
  \item \emph{Pages 89 et 90.} Groupes à Frobenius nilpotent et groupe
  formel associé à un groupe de BT (théorème~2.5) ; énoncé du théorème de
  déformation (4.1).
  \item \emph{Pages 91 à 95.} « Complexe cotangent relatif », en quatre
  parties (A) à (D), et un théorème sur les $\mathrm{Ext}$ au-dessus de
  $\mathbb{Z}/p^{n}$.
  \item \emph{Pages 96 à 107.} « Sorites sur groupes $p$-divisibles et
  groupes de B.T. » : les objets de Tate, d'abord dans une catégorie
  abélienne, puis parmi les schémas en groupes commutatifs, puis dans une
  suite exacte.
  \item \emph{Page 109.} Morphismes formellement lisses et un foncteur
  $\Phi$ que le dossier ne définit pas.
  \item \emph{Pages 112 à 114.} « Conférence à Nice » : le canevas d'un
  exposé, du groupe de Barsotti–Tate au foncteur de Dieudonné cristallin.
\end{enumerate}

Les pages 14, 78, 88 et 111 sont blanches. Les pages 98, 100, 102, 104, 108
et 110 sont des copies au duplicateur d'un texte étranger au sujet, au dos
desquelles il a écrit ; elles ne sont pas décrites. Les numéros au crayon
des archivistes sautent de 80 à 83 entre les pages 80 et 81 de la
numérisation : deux feuillets manquent peut-être, et la phrase sur laquelle
s'arrête la page 80 n'a pas de suite dans le dossier.

\subsection*{Le plan, et ce qui le remplit}

Le plan tapé de la page 77 donne la clé de l'ordre. Les étiquettes au crayon
des pages 76 (« I, II / Vecteurs de Witt et théorie de Dieudonné ») et 82
(« III Propriétés Générales des B.T. ») annoncent les chapitres I--II et III
de ce plan, et les pages qui les suivent les remplissent. Les pages 89 et
suivantes, numérotées 2, 4, 5, relèvent vraisemblablement des chapitres
IV--V (puissances divisées et cristaux, foncteur de Dieudonné : programme et
problèmes), sans que le dossier le dise. Le chapitre VI, « $F$-cristaux. Un
théorème de spécialisation », n'a aucune page.

Les trois chapitres de séminaire des pages 1 à 75 suivent ce même plan ; à
en juger par les titres que relèvent les notes de la transcription, le
chapitre III dactylographié reprend dans le même ordre les exemples, la
suite de composition et la proposition des pages 83 à 86\footnote{Ce
rapprochement est de nous, et ne repose que sur les intitulés notés aux
pages 67 à 75 de la transcription (schémas abéliens, groupes ind-étales,
tores, groupes toroïdaux, groupes de Lie formels et leurs conjectures,
courbe elliptique modulaire, suite de composition, « Proposition »
i)--iii)$^{\mathrm{bis}}$ et « Lemme » de factorisation). Qui a rédigé
quoi à partir de quoi, le dossier ne le dit pas.}. On lit donc côte à côte
une mise au propre et les brouillons dont elle procède, ou qui la
continuent.

\subsection*{Conventions}

\begin{itemize}
  \item $p$ est un nombre premier fixé « dans tout le cours » (page 79).
  « Groupe de BT » abrège groupe de Barsotti–Tate, c'est-à-dire groupe
  $p$-divisible au sens de Tate.
  \item $G(n) = \mathrm{Ker}(p^{n}\,\mathrm{id}_G)$, notation des pages,
  qu'on écrit aujourd'hui $G[p^{n}]$. La notation $G[i]$ est réservée,
  comme page 89, au noyau de la $i$-ième itérée du Frobenius relatif ; elle
  ne désigne jamais $G(i)$.
  \item $D^{*}$ est le foncteur de Dieudonné \emph{contravariant} ; sur un
  module de Dieudonné $M$, $F_M$ est $\sigma$-semi-linéaire et $V_M$
  $\sigma^{-1}$-semi-linéaire, $\sigma$ étant le Frobenius de $W = W(k)$ ;
  $F_M V_M = V_M F_M = p$. L'anneau de Dieudonné est
  $\mathcal{D} = W_{\sigma}[F, V]$.
  \item Pour un groupe plat $G$ sur $S$, $\ell^{G}$ est le complexe de
  co-Lie, en degrés homologiques $0$ et $1$ ; $\omega_G = H_0(\ell^{G})$,
  $n_G = H_1(\ell^{G})$, $t_G = \mathrm{Lie}(G)$, $\nu_G$.
  \item Un \emph{objet de Tate de degré $n$} est un couple $(A, f)$, $f$
  endomorphisme avec $f^{n} = 0$ et $\mathrm{Ker}\, f^{j} =
  \mathrm{Im}\, f^{n-j}$ pour tout $j$. Les pages disent tantôt
  « hauteur », tantôt « degré » pour cet $n$\footnote{Page 99 il remplace
  « hauteur » par « degré » ; pages 101 à 107 il revient à « hauteur ». On
  écrit « degré », pour ne pas confondre avec la hauteur d'un groupe de BT
  des pages 84 et 87.}.
\end{itemize}

\subsection*{Ce que le dossier annonce sans l'établir}

La suite de la page 80 ; les « Notes~I » dont la page 89 cite deux
propositions ; les conjectures de la page 83 sur les groupes de Lie
formels ; les démonstrations des théorèmes 2.5 et 4.1 ; la « lettre à
Illusie », citée deux fois et absente ; toute la page 93, écrite au
conditionnel ; les conditions (ii) et (ii bis) de la proposition 2.2,
laissées en blanc ; les démonstrations de la proposition (D) et du
théorème de la page 95 ; la « 1st conj. » de la page 106 ; la liste duale
de la page 107 ; le foncteur $\Phi$ de la page 109 ; et, dans le programme
de Nice, tout ce qui est introduit par « il se pourrait ».

\pagerange{1}{75}
\section*{I. Les tapuscrits du séminaire (pages 1 à 75)}

\pagerange{1}{27}
\subsection*{Chapitre I, « Préliminaires sur Witt » (pages 1 à 27)}

Deux frappes du même chapitre, sur deux machines : pages 1 à 13 et pages
15 à 27, la seconde avec de menues corrections de la première. La page de
titre porte au crayon, en capitales, « DELALE », d'une main
indéterminée\footnote{Le nom s'accorde avec le titre de l'inventaire, qui
le cite parmi quatre ; rien ne dit qu'il désigne l'auteur plutôt qu'un
destinataire.}. Selon les notes de la transcription, le chapitre traite
des schémas de vecteurs de Witt, de leurs morphismes remarquables, des
ind-schémas qu'on en tire, de l'anneau $W(k)$ pour $k$ parfait, puis du
Frobenius et de la Verschiebung d'un schéma en groupes, avec ses
propositions numérotées (2.3, 4.2) et une remarque. Aucune intervention
manuscrite.

\pagerange{28}{54}
\subsection*{Chapitre II, « Groupes finis localement libres et théorie classique de Dieudonné » (pages 28 à 54)}

Tapuscrit non signé, folioté à la main « II.1 » à « II.26 ». Les notes de
la transcription y relèvent : rappels sur les groupes finis localement
libres et la dualité de Cartier ; types particuliers de groupes et leurs
critères par le Frobenius et la Verschiebung ; décomposition canonique sur
un corps parfait ; théorie de Dieudonné sur un corps parfait (théorème~4.2
et ses corollaires 4.2.1, 4.3.3, le cas unipotent, le cas multiplicatif,
un lemme de décomposition) ; corollaires 5.1 et 5.2 et remarques 5.3 ;
une annexe construisant un quasi-inverse du foncteur de Dieudonné. Les
marges portent des corrections d'épreuve — croix, symboles repassés, mots
remplacés (« bigèbre », « directement », « bijectif ») — d'une main que le
feuillet ne permet pas d'identifier\footnote{Rien, dans ces corrections,
n'est propre à Grothendieck ; la transcription ne les lui attribue pas, et
on fait de même. Une inscription au crayon gommée, page 29, est
illisible.}.

\pagerange{55}{75}
\subsection*{Chapitre III, « Groupes de Barsotti–Tate » (pages 55 à 75)}

Tapuscrit non signé, folioté au crayon « III-1 » à « III-20 ». D'après les
notes : notations, platitude et critère de représentabilité (proposition
2.2, proposition 2.4) ; groupes de BT tronqués (définition 3.2) ; groupes
de BT (définitions 4.1, 4.2) ; « Sorites » (proposition 5.2 et sa
remarque, rang, systèmes $p$-adiques, dual, module de Dieudonné sur un
corps parfait) ; exemples (§6) ; suite de composition (§7, proposition 7.4
et un lemme). Mêmes corrections d'épreuve, de la même main indéterminée,
dont « Sorites » écrit sur « Suites »\footnote{Deux renvois du tapuscrit,
à SGA~3 VII et à EGA~IV, ont leur numéro laissé en blanc.}.

\pagerange{76}{77}
\section*{II. Le plan du cours (pages 76 et 77)}

La page 76 est une feuille de listing qui ne porte qu'une étiquette au
crayon : « I, II / Vecteurs de Witt et théorie de Dieudonné ». La page 77,
tapée sur sa machine, donne le plan :

\begin{enumerate}
  \item[I] Rappels sur les vecteurs de Witt, et
  \item[II] la théorie de Dieudonné.
  \item[III] Généralités sur les groupes de BT.
  \item[IV] Généralités sur les puissances divisées et les cristaux.
  \item[V] Foncteur de Dieudonné sur base à peu près quelconque :
  programme préliminaire et problèmes.
  \item[VI] $F$-cristaux. Un théorème de spécialisation pour les groupes
  de BT.
\end{enumerate}

C'est au crayon que la théorie de Dieudonné devient un chapitre à part
(« II »), que les chapitres suivants sont renumérotés et que le chapitre
V reçoit « et problèmes »\footnote{Tapé : « I Rappels sur les vecteurs de
Witt et la théorie de Dieudonné. » ; les numéros des trois dernières
lignes sont surchargés au crayon (la transcription lit « III » devenu
« IV », un numéro noirci remplacé par « V », un « V » complété en
« VI »). Le « à peu près quelconque » rend le « $\pm$ quelconque » de la
page.}. L'ajout dit l'état du travail : le
foncteur sur une base générale est un programme, et l'auteur le sait.

Le théorème de spécialisation du chapitre VI est, selon toute
vraisemblance, celui qu'on appelle aujourd'hui le théorème de
spécialisation de Grothendieck : le polygone de Newton d'une famille de
$F$-cristaux, ou de groupes de BT, ne peut que monter par
spécialisation\footnote{L'identification est de nous ; le dossier n'en
contient aucune page. L'énoncé est publié dans le volume tiré de ce
séminaire de Montréal (\emph{Groupes de Barsotti–Tate et cristaux de
Dieudonné}, Presses de l'Université de Montréal, 1974), cité de mémoire ;
Katz en a donné une démonstration (1979).}.

\pagerange{79}{81}
\section*{III. Vecteurs de Witt et théorie de Dieudonné (pages 79 à 81)}

\pagerange{79}{80}
\subsection*{Les vecteurs de Witt (pages 79 et 80)}

La page 79 s'intitule « Rappels sur la théorie de Dieudonné (sur corps
parfait $k$) » et renvoie, en marge, à Serre, \emph{Corps locaux}, et à
Demazure–Gabriel, \emph{Groupes algébriques}\footnote{La page écrit
« Schémas en Groupes » au-dessus du second nom, lu avec doute
« Demazure-Gabriel ».}.

Pour $n \geq 1$, $\underline{W}_n$ et $\underline{W} =
\varprojlim \underline{W}_n$ sont les schémas en anneaux des vecteurs de
Witt ($p$-typiques) sur $\mathbb{Z}$ ; comme schémas,
$\underline{W}_n = \mathbb{A}^{n}$ et $\underline{W} = \mathbb{A}^{\mathbb{N}}$.
L'addition et la multiplication sont les seules pour lesquelles les
composantes fantômes
\[
\phi_i(x_1, \dots, x_n) = x_1^{p^{i-1}} + p\, x_2^{p^{i-2}} + \dots
+ p^{i-1} x_i \qquad (1 \leq i \leq n)
\]
définissent un homomorphisme d'anneaux
$\Phi = (\phi_1, \dots, \phi_n) \colon \underline{W}_n \to
\mathbb{G}_a^{n}$\footnote{La page écrit $\underline{O}^{n}$ pour
l'anneau $\mathbb{G}_a^{n}$ et $\underline{E}^{n}$ pour l'espace affine
sous-jacent. Sur $\mathbb{Z}$, $\Phi$ n'est pas un isomorphisme ; il le
devient après inversion de $p$, et c'est ce qui caractérise les lois.}.
On dispose de
\begin{itemize}
  \item la \emph{restriction} $R \colon \underline{W}_{n+1} \to
  \underline{W}_n$, $(x_1, \dots, x_{n+1}) \mapsto (x_1, \dots, x_n)$,
  morphisme d'anneaux ;
  \item la \emph{Verschiebung} $V \colon \underline{W}_n \to
  \underline{W}_n$, $(x_1, \dots, x_n) \mapsto (0, x_1, \dots,
  x_{n-1})$, additive, et par passage à la limite $V \colon \underline{W}
  \to \underline{W}$ ;
  \item le décalage $T \colon \underline{W}_n \to \underline{W}_{n+1}$,
  $(x_1, \dots, x_n) \mapsto (0, x_1, \dots, x_n)$, additif,
\end{itemize}
reliés par $R_n T_n = V_n$ et $T_n R_n = V_{n+1}$. Enfin le morphisme de
schémas $F \colon (x_1, \dots, x_n) \mapsto (x_1^{p}, \dots, x_n^{p})$ est,
\emph{en caractéristique $p$}, un endomorphisme d'anneaux de
$\underline{W}_{n,\mathbb{F}_p}$ et de $\underline{W}_{\mathbb{F}_p}$ :
c'est le Frobenius, et l'on a
\[
FV = VF = p \cdot \mathrm{id} \qquad \text{en caractéristique } p
\]
\footnote{Sur $\mathbb{Z}$, l'application $(x_i) \mapsto (x_i^{p})$ n'est
qu'un morphisme de schémas : ce n'est pas un morphisme d'anneaux. Le
Frobenius des vecteurs de Witt sur $\mathbb{Z}$ est un autre morphisme
polynomial, $\underline{W}_{n+1} \to \underline{W}_n$, caractérisé par
$\phi_i \circ F = \phi_{i+1}$, qui se réduit modulo $p$ sur celui-ci. La
page définit $F$ sur $\underline{W}_n$ sans préciser, puis dit « en
caractéristique $p$ » ; la précision est de nous. La page a d'abord tapé,
puis barré, « $V$, $D$ sont des morphismes additifs ».}.

On introduit l'ind-schéma
$\underline{W}_{\to} = \varinjlim (\underline{W}_n, T)$, limite inductive
le long des décalages, qui comme ind-schéma est
$\mathbb{A}^{(\mathbb{N})}$\footnote{La page dit « schéma en groupes » ;
c'est un ind-schéma en groupes (un faisceau en groupes limite inductive de
schémas). Elle nomme les transitions « du type $D$ » : la lettre $T$ de la
page 79 est retapée sur un $D$.}. En caractéristique $p$, le diagramme
suivant commute :

\begin{tikzcd}
\underline{W}_n \arrow[r, "T"] \arrow[d, "F^{n-1}"'] & \underline{W}_{n+1} \arrow[d, "F^{n}"] \\
\underline{W}_n \arrow[r, "p"] & \underline{W}_{n+1}
\end{tikzcd}

où la flèche du bas est induite par la multiplication par $p$, qui
passe au quotient $\underline{W}_n \to \underline{W}_{n+1}$ et vaut
$T \circ F$\footnote{Que $p = T \circ F$ comme application
$\underline{W}_n \to \underline{W}_{n+1}$ en caractéristique $p$ découle
de $p = VF$ ; c'est exactement ce que dit le premier carré de la page, à
flèches verticales $F$ et $\mathrm{id}$, dont la flèche du bas, manuscrite,
est lue « $r$ » avec doute. Nous la lisons $p$, seule lecture qui fasse
commuter le carré ; la lecture est de nous.}. D'où un homomorphisme
\[
u \colon \underline{W}_{\to,\mathbb{F}_p} \longrightarrow
\varinjlim\, (\underline{W}_{n,\mathbb{F}_p}, p),
\]
égal à $F^{n-1}$ sur le $n$-ième terme, qui est un isomorphisme sur les
points à valeurs dans un anneau parfait, puisque $F$ y est
bijectif\footnote{La page a d'abord tapé « corps », puis « anneau »
au-dessus.}.

Soit $k$ un corps parfait de caractéristique $p$. Alors
$W = \underline{W}(k) = \varprojlim W_n(k)$ est un anneau de valuation
discrète complet, de corps résiduel $W_1(k) = k$, d'idéal maximal $pW$, de
corps des fractions $K = W[1/p]$ de caractéristique nulle, et $u$ induit
un isomorphisme de groupes
\[
u \colon \underline{W}_{\to}(k) \xrightarrow{\ \sim\ } K/W ,
\]
le $n$-ième terme $W_n(k) \simeq W/p^{n}W$ s'envoyant sur $p^{-n}W/W$ par
$x \mapsto F^{n-1}(x)/p^{n}$. La page 80 s'arrête au milieu de la phrase
suivante (« On va introduire sur $\underline{W}_k$ une structure de schéma
en groupes avec »), et la suite manque\footnote{Le feuillet suivant de la
numérisation, manuscrit, n'en est pas la suite, et la numérotation au
crayon saute de deux numéros à cet endroit.}.

\pagerange{81}{81}
\subsection*{La décomposition sur un corps parfait et le théorème de Dieudonné (page 81)}

La page 81, manuscrite, commence en cours d'exposé ; son « ceci » renvoie à
une page qui n'est pas dans le dossier. Soit $k$ un corps parfait de
caractéristique $p$ et $G$ un $p$-groupe fini commutatif sur $k$
(schéma en groupes fini, commutatif, d'ordre une puissance de $p$). Il se
décompose canoniquement en
\[
G \simeq G_{\mathrm{et}} \times G_{\mathrm{mult}} \times G_{\mathrm{bi}},
\]
partie étale, partie de type multiplicatif et partie bi-infinitésimale
(infinitésimale de dual infinitésimal), et la catégorie des $p$-groupes
finis est équivalente au produit des trois catégories
correspondantes\footnote{Nom moderne : la décomposition en parties
étale-locale, locale-étale et locale-locale ; pour un $p$-groupe la
partie étale-étale est nulle. La page écrit « somme directe ».}. Les deux
premières sont connues par la théorie de Galois (et la dualité de Cartier
pour la seconde) ; la troisième demande la théorie de Dieudonné. Le
Frobenius et la Verschiebung lisent la décomposition :
\[
\begin{array}{lcl}
F_G \text{ nilpotent} & \iff & G_{\mathrm{et}} = 0 \iff G^{*} \text{ unipotent}, \\
V_G \text{ nilpotent} & \iff & G_{\mathrm{mult}} = 0 \iff G \text{ unipotent}, \\
F_G \text{ isomorphisme} & \iff & G \text{ étale}, \\
V_G \text{ isomorphisme} & \iff & G \text{ de type multiplicatif},
\end{array}
\]
où « nilpotent » s'entend des itérés $F^{n}_{G/k} \colon G \to G^{(p^{n})}$
et $V^{n}_{G/k}$\footnote{C'est la note marginale de la page, reliée aux
deux « nilpotent » cerclés : « préciser ce que ça veut dire, en
introduisant Frob.\ et Versch.\ \emph{itérés} ».}.

\emph{Théorème de Dieudonné.} Le foncteur $G \mapsto D^{*}(G)$ est une
anti-équivalence entre la catégorie des $p$-groupes finis commutatifs sur
$k$ et celle des modules de Dieudonné de longueur finie : $W$-modules $M$
de longueur finie munis de $F_M, V_M \colon M \to M$ tels que
\[
F_M \lambda = \lambda^{\sigma} F_M, \qquad
\lambda V_M = V_M \lambda^{\sigma}, \qquad
F_M V_M = V_M F_M = p \cdot \mathrm{id}
\qquad (\lambda \in W)
\]
\footnote{La page écrit deux fois « $FV = p\cdot\mathrm{id}$ » ; la
seconde est pour $VF$. Elle ajoute « anti » au-dessus de « équivalence »
dans l'interligne.}. Autrement dit $F_M$ est une application linéaire
$M^{(\sigma)} = W \otimes_{\sigma, W} M \to M$ et $V_M$ une application
linéaire $M \to M^{(\sigma)}$. Comme $D^{*}$ commute au changement de
base, $D^{*}(G^{(p)}) \simeq D^{*}(G)^{(\sigma)}$, et le théorème précise
\[
F_{D^{*}(G)} = D^{*}(F_G), \qquad V_{D^{*}(G)} = D^{*}(V_G),
\]
d'où le dictionnaire
\[
\begin{array}{lcl}
G \text{ unipotent} & \iff & V_M \text{ nilpotent}, \\
G^{*} \text{ unipotent (i.e. } G \text{ infinitésimal)} & \iff & F_M \text{ nilpotent}, \\
G \text{ étale} & \iff & F_M \text{ bijectif}, \\
G^{*} \text{ étale (i.e. } G \text{ de type multiplicatif)} & \iff & V_M \text{ bijectif}.
\end{array}
\]
Sur les trois groupes d'ordre $p$, on contrôle : pour $\mathbb{Z}/p$,
$M = k$ avec $F$ bijectif et $V = 0$ ; pour $\mu_p$, $F = 0$ et $V$
bijectif ; pour $\alpha_p$, $F = V = 0$\footnote{Vérification de nous,
avec la convention contravariante.}. Le chapitre de séminaire des pages 28
à 54 porte, d'après ses notes, le même théorème sous le numéro 4.2.

\pagerange{82}{87}
\section*{IV. Propriétés générales des groupes de BT (pages 82 à 87)}

La page 82 est une étiquette au crayon, « III Propriétés Générales des
B.T. » ; les pages 83 à 87 sont tapées, avec des ajouts à l'encre et au
crayon. Un groupe de BT sur $S$ est un faisceau fppf en groupes
commutatifs $G$, $p$-divisible ($p\,\mathrm{id}_G$ épimorphique), de
$p$-torsion ($G = \varinjlim G(n)$), avec $G(1)$ fini localement libre ;
c'est la définition qu'énonce la page 112, et celle des tapuscrits.

\pagerange{83}{84}
\subsection*{Exemples (pages 83 et 84)}

\emph{a) Schémas abéliens et tores.} Soit $G$ un schéma en groupes lisse
sur $S$, extension d'un schéma abélien $A$ par un tore $T$. Alors
$p\,\mathrm{id}_G$ est un épimorphisme et les $G(n)$ sont finis
localement libres (on commence par le schéma abélien) ; donc
$G(\infty) = \varinjlim G(n)$ est un groupe de BT, de hauteur $2a + m$
si $A$ est de dimension relative $a$ et $T$ de dimension relative
$m$\footnote{La hauteur est l'ajout au crayon de la page, dans une bulle.}.
Cas important : $\mu = \mathbb{G}_m(\infty) = \varinjlim \mu_{p^{n}}$.

\emph{b) Groupes de Lie formels.} Supposons $p$ localement nilpotent sur
$S$ et soit $G$ un groupe de Lie formel (commutatif) sur $S$ ; alors
$G = \varinjlim G(n)$. Supposons $p\,\mathrm{id}_G$ épimorphique.
Lorsque $S$ est le spectre d'un anneau local artinien de point fermé $s$,
cela revient à demander que $p\,\mathrm{id}_{G_s}$ soit un épimorphisme,
ou encore, selon la page, que $G_{\bar s}$ ne contienne pas de sous-groupe
formel isomorphe à $\widehat{\mathbb{G}}_a$\footnote{Ce dernier critère
est énoncé sans démonstration ; nous le laissons comme sien. Sur un corps,
l'épimorphie de $p$ revient à ce que le groupe formel soit de hauteur
finie.}. Sur une base qui n'est pas un point, un trait par exemple, être
de BT n'est plus une condition fibre à fibre : le groupe formel associé à
la courbe elliptique modulaire, au voisinage d'un point d'invariant de
Hasse nul, a des fibres de hauteur $1$ en général et de hauteur $2$ en ce
point, de sorte que $G(1)$ ne peut y être fini localement libre, son rang
sautant de $p$ à $p^{2}$\footnote{L'explication par le saut de rang est de
nous ; la page renvoie à une note (*) au crayon : « La condition suppl.\ à
rajouter devrait être : fibres de "hauteur" localement constante. »}.

La page conjecture alors que, si $p\,\mathrm{id}_G$ est un épimorphisme,
$G(1)$ et donc les $G(n)$ sont finis localement libres — et plus
généralement que le noyau d'un épimorphisme de groupes de Lie formels de
même dimension relative est fini localement libre, avec la réciproque.
C'est vrai si $S$ est artinien, et on l'\emph{ajoute} comme hypothèse
dans le cas général\footnote{La page tape « les $G(1)$ dont les $G(n)$
sont finis localement libres », qu'on lit « donc ». Nous ne savons pas
dire, de mémoire, si la conjecture est démontrée sous cette forme ;
l'usage moderne la contourne en définissant un groupe formel
$p$-divisible par la condition que $p\,\mathrm{id}_G$ soit une isogénie,
c'est-à-dire un épimorphisme de noyau fini localement libre.}. Sous ces
hypothèses $G$ est un groupe de BT, et les $G(n)$ sont radiciels. La page
annonce, pour plus tard, une équivalence de catégories :
\[
\left\{ \begin{array}{c} \text{groupes de BT sur } S \\ \text{à fibres connexes} \end{array} \right\}
\simeq
\left\{ \begin{array}{c} \text{groupes de Lie formels } G \text{ sur } S, \\
p\,\mathrm{id}_G \text{ épimorphique}, \ G(1) \text{ fini loc.\ libre} \end{array} \right\},
\]
« à fibres connexes » signifiant que $G(1)$, ou de façon équivalente
chaque $G(n)$, est radiciel\footnote{Énoncé moderne : si $p$ est
localement nilpotent sur $S$, les groupes $p$-divisibles connexes sont
les groupes formels $p$-divisibles ; c'est un théorème de Tate sur un
anneau local noethérien complet (Driebergen, 1966), et de Messing sur une
base où $p$ est localement nilpotent (\emph{The crystals associated to
Barsotti–Tate groups}, 1972). La page ne démontre pas l'équivalence.}.
Le rang du groupe de BT ainsi obtenu s'appelle la \emph{hauteur} du
groupe formel.

\emph{c) Groupes de BT ind-étales.} Le foncteur $G \mapsto T_p(G) =
\varprojlim G(n)$ est une équivalence entre groupes de BT ind-étales et
faisceaux $p$-adiques étales lisses sans torsion (« constants tordus ») ;
si $S$ est connexe et pointé par $\bar s$, ceux-ci sont les représentations
continues de $\pi_1(S, \bar s)$ dans des $\mathbb{Z}_p$-modules libres de
type fini.

\emph{Groupes de BT toroïdaux} : ceux dont $G(1)$, donc chaque $G(n)$, est
de type multiplicatif. Le système $(G(n)^{*})_n$ des duaux de Cartier
donne une anti-équivalence avec les faisceaux $p$-adiques lisses sans
torsion ; composée avec la dualité de ces faisceaux, elle devient une
équivalence, qui s'écrit
\[
G \longmapsto \underline{\mathrm{Hom}}_{\mathrm{gr}}(\mu, G)
= \bigl( \underline{\mathrm{Hom}}_{\mathrm{gr}}(\mu_{p^{n}}, G(n)) \bigr)_{n \geq 1}
\]
(abus de notation que la page signale). Si les caractéristiques
résiduelles de $S$ sont égales à $p$, les groupes de BT toroïdaux sont à
fibres connexes ; et réciproquement, s'ils le sont tous, les
caractéristiques résiduelles valent $p$\footnote{La page écrit « (et
réciproquement) », et une note marginale reliée à ces mots précise
« (s'il en est ainsi, car.\ rés.\ sont $p$) ». Nous lisons la réciproque
ainsi ; lue comme « un groupe de BT à fibres connexes est toroïdal », elle
serait fausse ($\widehat{E}(\infty)$ pour une courbe elliptique
supersingulière est connexe et n'est pas toroïdal). Exemple de nous.}.

\pagerange{84}{86}
\subsection*{La filtration connexe-étale, et quand elle existe (pages 84 à 86)}

\emph{d) Groupes de BT comme extensions.} Supposons d'abord $S$ réduit à
un point $s$ (spectre d'un anneau local artinien, par exemple). Pour tout
groupe fini $H$ sur $S$, la composante neutre $H^{\circ}$ est un
sous-groupe, fonctoriel en $H$ ; si $H$ est fini localement libre,
$H^{\circ}$ l'est aussi, et $H/H^{\circ}$ est fini localement libre et
étale. Ainsi $H$ est, canoniquement, extension d'un groupe étale par un
groupe infinitésimal. Si $k(s)$ est de caractéristique $p$, la dualité de
Cartier fournit de même un plus grand sous-groupe de type multiplicatif
$H^{\mathrm{tm}} \subset H^{\circ}$, d'où une filtration à trois crans
$0 \subset H^{\mathrm{tm}} \subset H^{\circ} \subset H$, « qu'on avait eu
le tort de ne définir précédemment que sur un \emph{corps} de base ».

Par passage à la limite, tout groupe de BT $G$ sur un tel $S$ a une
filtration canonique et fonctorielle par des sous-groupes de BT
\[
0 \subset G^{\mathrm{tm}} = \varinjlim G(n)^{\mathrm{tm}} \subset
G^{\circ} = \varinjlim G(n)^{\circ} \subset G,
\]
$G^{\circ}$ plus grand sous-groupe de BT à fibres connexes et (si
$\mathrm{car}\, k(s) = p$) $G^{\mathrm{tm}}$ plus grand sous-groupe de BT
toroïdal. En particulier $G$ est extension du groupe de BT ind-étale
$G/G^{\circ}$ par le groupe de BT connexe $G^{\circ}$, lequel, si
$\mathrm{car}\, k(s) = p$ et que l'idéal maximal de $\mathcal{O}_{S,s}$
est nilpotent, est un groupe de Lie formel (b) ; et $G$ est une extension
triple : ind-étale, par connexe à dual connexe, par toroïdal. C'est la
suite connexe-étale d'aujourd'hui.

Sur une base qui n'est pas un point — un trait de caractéristique $p$ —,
un groupe de BT n'est plus en général extension d'un ind-étale par un
connexe : $A(\infty)$, pour la courbe elliptique modulaire, échoue encore
aux points d'invariant de Hasse nul, où la partie étale disparaît.

\emph{Proposition.} Pour un groupe de BT $G$ sur $S$, conditions
équivalentes :
\begin{itemize}
  \item[(i)] $G$ est extension d'un groupe de BT ind-étale par un groupe
  de BT à fibres connexes ;
  \item[(ii)] pour tout $n$, $G(n)$ est extension d'un groupe fini étale
  par un groupe fini radiciel ;
  \item[(iii)] de même pour $n = 1$ ;
  \item[(iv)] la fonction $s \mapsto$ rang séparable de $G(1)_s$, nombre de
  points géométriques $\mathrm{card}\, G(1)(\bar s)$, est localement
  constante\footnote{La page écrit $\mathrm{rang\ sép}\, G_s =
  \mathrm{card}\, G(\bar s)$, qui serait infini ; il s'agit de $G(1)$.
  Ce rang vaut $p^{h_{\mathrm{et}}(s)}$, et (iv) dit que la hauteur étale
  est localement constante.}.
\end{itemize}
Elle résulte, dit la page, du lemme suivant, « qui figure dans EGA~IV ? ».

\emph{Lemme.} Soit $X \to S$ fini localement libre. Il se factorise en
$X \to Y \to S$, avec $X \to Y$ radiciel surjectif (fini localement libre)
et $Y \to S$ fini étale, si et seulement si la fonction
$s \mapsto \mathrm{card}\, X(\bar s)$ est localement constante. La
factorisation est alors unique et fonctorielle en $X$ ; si $X$ est un
schéma en groupes, elle est une structure d'extension, et $X \mapsto
X^{\mathrm{et}}$, $X \mapsto X^{\circ}$ sont exacts sur les suites exactes
courtes\footnote{La page tape « se factorise en un morphisme étale suivi
d'un morphisme radiciel surjectif », après avoir barré « radiciel » ;
l'ordre est l'inverse de celui qui est vrai, et que donne le lemme du
tapuscrit de la page 75 selon sa note : d'abord le radiciel, puis l'étale.
Elle écrit aussi $G$ pour $X$ dans la fin de l'énoncé. Nous n'avons pas
vérifié le renvoi à EGA~IV.}. Le reste de la page 86 est blanc.

\pagerange{87}{87}
\subsection*{Extensions, conoyaux, hauteur (page 87)}

\emph{Proposition.} (i) Toute extension d'un groupe de BT par un groupe
de BT est un groupe de BT. (ii) Le conoyau de tout monomorphisme de
groupes de BT est un groupe de BT.

La page barre un (iii) — « le noyau de tout épimorphisme est de BT » —,
et à juste titre : le noyau de $p\,\mathrm{id}_G$ est $G(1)$, qui n'est
pas de BT\footnote{L'exemple est de nous ; le tapuscrit du chapitre III,
page 64, porte d'après sa note une remarque dans ce sens.}.

\emph{Démonstration.} Soit $0 \to G' \to G \to G'' \to 0$ exacte. Le
lemme du serpent pour la multiplication par $p$ donne la suite exacte
\[
0 \to G'(1) \to G(1) \to G''(1) \to G'/pG' \to G/pG \to G''/pG'' \to 0 .
\]
Dans le cas (i), $G'/pG' = G''/pG'' = 0$, donc $G/pG = 0$ ; dans le cas
(ii), $G'/pG' = G/pG = 0$, donc $G''/pG'' = 0$. Dans les deux cas
$0 \to G'(1) \to G(1) \to G''(1) \to 0$ est exacte, d'où $G(1)$ (cas i,
extension de deux groupes finis localement libres) ou $G''(1)$ (cas ii,
quotient d'un fini localement libre par un sous-groupe fini localement
libre) est fini localement libre. Les autres conditions sont claires.

\emph{Hauteur.} Le rang de $G(1)$ est en chaque point de la forme $p^{r}$,
$r$ fonction localement constante sur $S$ ; $r$ est le \emph{rang}, ou
\emph{hauteur}, du groupe de BT, et $G(n)$ est de rang $p^{rn}$.

\pagerange{89}{90}
\section*{V. Frobenius nilpotent, groupe formel associé, déformations (pages 89 et 90)}

\pagerange{89}{89}
\subsection*{Groupes à Frobenius nilpotent (page 89)}

« 2. Sur certains schémas en groupes finis loc.\ libres en car.\ $p$
(Applications aux BT) ». Ici $S$ est de caractéristique $p$, et la page
note $\underline{f}$ un endomorphisme qu'elle ne définit pas ; nous le
lisons comme le Frobenius relatif $F$, avec
$G[i] = \mathrm{Ker}(F^{i}_{G/S} \colon G \to G^{(p^{i})})$\footnote{La
lecture est de nous ; elle s'appuie sur la suite de la page (« lisse à
l'ordre $p^{n} - 1$ », $\mathrm{Inf}^{k}$, groupes de Lie formels), qui
n'a de sens que pour le Frobenius. Les propositions 2.1 et 2.2 renvoient
aux propositions 14 et 1 de « Notes~I », qui ne sont pas dans le
dossier.}. On suppose $G = G[n]$, et l'on fixe $1 \leq i \leq n-1$.

\emph{Proposition 2.2.} Conditions équivalentes : (i) $F^{i}$, qui envoie
$G$ dans $G^{(p^{i})}[n-i] = G[n-i]^{(p^{i})}$, est plat sur son but,
donc fidèlement plat ; (i bis) il est couvrant ; (ii), (ii bis) deux
formulations par le gradué associé, l'une faisceautique, l'autre en
schémas en groupes, que la page laisse en blanc\footnote{La page écrit
le but $G[n-i]^{p^{n-i}}$ ; l'exposant doit être $p^{i}$, puisque $F^{i}$
va dans $G^{(p^{i})}$. Une condition « (ii) Les groupes $G[i]$ sont tous
plats » est tapée puis surfrappée.}. C'est, au Frobenius près, la notion
d'objet de Tate des pages 97 à 107 (« $f^{n-j} \colon G \to
\mathrm{Ker}\, f^{j}$ fidèlement plat »), ici avec $f = F$ et la torsion
par $(p)$\footnote{Le rapprochement est de nous.}.

Un tel groupe est lisse à l'ordre $p^{n} - 1$. Combinant 2.1 et 2.2 :

\emph{Corollaire 2.3.} Si $G$ est un groupe de BT sur $S$ de
caractéristique $p$, alors pour tout $n > 0$, $G(n)$ est lisse à l'ordre
$p^{n} - 1$ sur $S$, et $\mathrm{Inf}^{k}(G) \subset G[n]$ pour
$k \leq p^{n} - 1$, où $G[n] = G(n)[n] = G(i)[n]$ pour $i \geq
n$\footnote{Note au crayon de la page. L'inclusion
$\mathrm{Inf}^{k} \subset \mathrm{Ker}\, F^{n}$ pour $k + 1 \leq p^{n}$
est immédiate : $F^{n}$ élève les fonctions à la puissance $p^{n}$.}.

\emph{Proposition 2.4.} La catégorie des groupes de Lie formels
commutatifs sur $S$ de caractéristique $p$ est équivalente à celle des
systèmes inductifs $F$-coadiques $(G[n])_n$ de schémas en groupes finis
localement libres satisfaisant aux conditions de 2.2.

\emph{Théorème 2.5.} Soit $G$ un groupe de BT sur $S$ de caractéristique
$p$. Alors
\[
\overline{G} = \mathrm{Inf}^{\infty}(G) = \varinjlim_n G(n)[n]
\]
est un groupe de Lie formel sur $S$, avec $\overline{G}[n] = G(n)[n]$,
et pour $k \leq p^{n} - 1$,
\[
\mathrm{Inf}^{k} G = \mathrm{Inf}^{k} \overline{G} = \mathrm{Inf}^{k} G(n)
\subset G(n)[n] \subset G(n).
\]
C'est le groupe formel $\widehat{G}$ associé à un groupe de BT. Une note
en marge demande d'en « donner [la] relation $d + d^{*} = h$ » : la
dimension de $\widehat{G}$ plus celle de $\widehat{G^{*}}$ égale la
hauteur\footnote{Relation classique ; la page ne l'établit pas.}. Les
démonstrations ne sont pas dans le dossier.

\pagerange{90}{90}
\subsection*{Le théorème de déformation (page 90)}

« 4. Déformations de groupes de BT : énoncés. » Le résultat, « qui sera
essentiel pour la théorie de Dieudonné », renvoie à la « lettre Illusie,
6.4 ».

\emph{Théorème 4.1.} Soit $S_0 \hookrightarrow S$ une immersion fermée
définie par un idéal nilpotent $\mathcal{I}$, $S$ affine et $p$
localement nilpotent sur $S$, et $G_0$ un groupe de BT sur
$S_0$\footnote{La page ne suppose $p$ nilpotent qu'au c) ($p^{N} 1_{S_0}
= 0$, qui l'entraîne sur $S$) ; nous l'ajoutons pour a) et b), comme le
fait l'énoncé publié par Illusie (« Déformations de groupes de
Barsotti–Tate, d'après A.~Grothendieck », Astérisque~127, 1985), cité de
mémoire.}. Notons $E(G_0, S)$ l'ensemble des classes de groupes de BT sur
$S$ qui prolongent $G_0$, et de même $E(G_0(n), S)$ pour les groupes de BT
tronqués d'échelon $n$. Alors :
\begin{itemize}
  \item[a)] il existe un groupe de BT $G$ sur $S$ prolongeant $G_0$ ;
  \item[b)] l'application $E(G_0, S) \to E(G_0(n), S)$, $G \mapsto G(n)$,
  est surjective ;
  \item[c)] si $\mathcal{I}^{k} = 0$ et $p^{N} 1_{S_0} = 0$, elle est
  bijective pour $n \geq kN$ ;
  \item[d)] si $\mathcal{I}^{2} = 0$, $E(G_0, S)$ est un torseur sous
  $\Gamma(S_0, t_{G_0^{*}} \otimes t_{G_0} \otimes \mathcal{I})$\footnote{La
  page écrit « nilpotente d'ordre $k$ » et, au d), « $k = 1$ » ; nous
  lisons $\mathcal{I}^{k} = 0$ et le cas d'un idéal de carré nul. Elle
  écrit le torseur sous $\underline{t}_{G_0^{*}} \otimes
  \underline{t}_{G_0}$ sans $\mathcal{I}$, qui doit y figurer. Nous ne
  vérifions pas la borne $kN$.}.
\end{itemize}
Le résultat « s'obtient via un résultat de déformation pour les groupes de
BT tronqués (je ne sais pas procéder directement !) ». Les pages 94 et 95
en donnent les ingrédients. Le n°~5 annoncé ensuite, « Rapport sur le
complexe cotangent relatif », est la suite manuscrite.

\pagerange{91}{95}
\section*{VI. Le complexe cotangent relatif (pages 91 à 95)}

Quatre parties, marquées (A) à (D), qui forment le « rapport » annoncé à la
page 90 : d'abord les généralités, puis les groupes plats, puis les
déformations, puis les groupes de BT tronqués.

\pagerange{91}{91}
\subsection*{(A) Généralités (page 91)}

Pour $X \to Y$, le complexe cotangent $L^{X/Y}$ est le « dérivé » du
foncteur $\Omega^{1}_{X/Y}$, dont on admet la définition (Illusie) ; il
est fonctoriel. On calcule son tronqué $\tau_{\leq 1} L^{X/Y}$ par un
plongement dans un schéma (formellement) lisse ; il est quasi-cohérent,
cohérent sous des hypothèses de finitude ; si $X \to Y$ est d'intersection
complète relative, $H_i(L^{X/Y}) = 0$ pour $i \geq 2$. Il a un triangle
distingué de transitivité, d'où une suite exacte longue, et commute au
changement de base $S' \to S$ lorsque
$\mathrm{Tor}_i^{\mathcal{O}_S}(\mathcal{O}_X, \mathcal{O}_{S'}) = 0$
pour $i > 0$ (par exemple si $X \to S$ ou $S' \to S$ est plat).

Trois applications, qu'on donne dans l'ordre logique\footnote{La page
les numérote 2, 1, 3 dans cet ordre de haut en bas.}.

\emph{1. Extensions d'algèbres.} Les classes d'extensions de la
$\mathcal{O}_S$-algèbre $\mathcal{O}_X$ par un $\mathcal{O}_X$-module
$\mathcal{J}$ de carré nul forment $\mathrm{Ext}^{1}_{\mathcal{O}_X}(L^{X/S},
\mathcal{J})$, et les automorphismes d'une extension
$\mathrm{Ext}^{0}_{\mathcal{O}_X}(L^{X/S}, \mathcal{J}) =
\mathrm{Hom}(\Omega^{1}_{X/S}, \mathcal{J})$, les dérivations.

\emph{2. Prolongement de morphismes.} Soit $Y_0 = V(\mathcal{J}) \subset
Y$ avec $\mathcal{J}^{2} = 0$, au-dessus de $S$, et $f_0 \colon Y_0 \to
X$ un $S$-morphisme. L'obstruction à prolonger $f_0$ en $f \colon Y \to X$
est dans $\mathrm{Ext}^{1}_{\mathcal{O}_{Y_0}}(Lf_0^{*} L^{X/S},
\mathcal{J})$ ; si elle s'annule, les prolongements forment un torseur
sous $\mathrm{Ext}^{0}_{\mathcal{O}_{Y_0}}(Lf_0^{*} L^{X/S},
\mathcal{J})$.

\emph{3. Déformations plates.} Soit $S_0 = V(\mathcal{J}) \subset S$ avec
$\mathcal{J}^{2} = 0$ et $X_0$ plat sur $S_0$. Posons
$\mathcal{J}_{X_0} = \mathcal{J} \otimes_{\mathcal{O}_{S_0}}
\mathcal{O}_{X_0}$. L'obstruction à déformer $X_0$ en un $X$ plat sur $S$
est dans $\mathrm{Ext}^{2}_{\mathcal{O}_{X_0}}(L^{X_0/S_0},
\mathcal{J}_{X_0})$ ; les déformations forment, si elle s'annule, un
torseur sous $\mathrm{Ext}^{1}$, et les automorphismes d'une déformation
sont $\mathrm{Ext}^{0}$. C'est la théorie d'Illusie (\emph{Complexe
cotangent et déformations}~I, 1971)\footnote{Cité de mémoire. La page
ajoute en marge « Lier avec Ex.\,? ».}.

\pagerange{92}{93}
\subsection*{(B) Groupes plats : le complexe de co-Lie (pages 92 et 93)}

Soit $G$ un schéma en groupes commutatif plat sur $S$, de section unité
$e$. On pose
\[
\ell^{G} = Le^{*} L^{G/S} \quad \text{(complexe de co-Lie)}, \qquad
\check{\ell}^{G} = R\underline{\mathrm{Hom}}(\ell^{G}, \mathcal{O}_S)
\quad \text{(complexe de Lie)},
\]
notion que la page attribue à Mazur et Roberts\footnote{Note marginale.
Nous ne vérifions pas l'attribution. La page signale aussi, d'un mot,
les « variances » : $\ell$ et $\check{\ell}$ sont duales.}. Si $G$ est
localement de présentation finie, il est d'intersection complète relative,
et $\ell^{G}$ est parfait d'amplitude parfaite dans $[-1, 0]$
(cohomologiquement), $\check{\ell}^{G}$ dans $[0, 1]$. On note
\[
\omega_G = H_0(\ell^{G}), \quad n_G = H_1(\ell^{G}), \quad
t_G = H^{0}(\check{\ell}^{G}), \quad \nu_G = H^{1}(\check{\ell}^{G}),
\]
de sorte que $t_G = \check{\omega}_G$ et $n_G = \check{\nu}_G$, et le rang
de $\ell^{G}$, qui vaut $\mathrm{rg}\,\omega_G - \mathrm{rg}\, n_G$ quand
ces deux modules sont localement libres, est la dimension relative de
$G$.

Une suite exacte $0 \to G' \to G \to G'' \to 0$ de groupes plats donne un
triangle distingué $\ell^{G''} \to \ell^{G} \to \ell^{G'} \to
\ell^{G''}[1]$, d'où, si $G'$ est localement de présentation finie, la
suite exacte
\[
0 \to n_{G''} \to n_G \to n_{G'} \to \omega_{G''} \to \omega_G \to
\omega_{G'} \to 0
\]
\footnote{La page dit « suite exacte courte » ; les deux premiers indices
sont surchargés, et l'ordre est rétabli d'après celui des $\omega$.}.

\emph{Formule remarquable.} Pour $G$ fini localement libre, de dual de
Cartier $G^{*}$, et $\mathcal{J}$ un $\mathcal{O}_S$-module
quasi-cohérent,
\[
R\underline{\mathrm{Hom}}_{\mathcal{O}_S}(\ell^{G}, \mathcal{J}) \simeq
\tau_{\leq 1} R\underline{\mathrm{Hom}}(G^{*}, \mathcal{J}),
\]
le second $R\underline{\mathrm{Hom}}$ étant pris entre faisceaux
abéliens. En degré $0$ : $t_G \simeq \underline{\mathrm{Hom}}(G^{*},
\mathbb{G}_a)$ ; en degré $1$ : $\nu_G \simeq
\underline{\mathrm{Ext}}^{1}(G^{*}, \mathbb{G}_a)$\footnote{La page
n'énonce pas d'hypothèse sur $G$ ; la dualité $G \mapsto G^{*}$ que nomme
sa note marginale demande $G$ fini localement libre (ou de BT), et nous la
supposons. La note marginale écrit $\underline{\mathrm{Hom}}(G^{*},
\mathcal{O}_S) \simeq \underline{t}_{G^{*}}$, ce qui contredit la formule
encadrée ; celle-ci donne $t_G$, et c'est ce qui est juste
($\mathrm{Lie}(\mu_p) = \mathbb{G}_a = \underline{\mathrm{Hom}}(\mathbb{Z}/p,
\mathbb{G}_a)$ en caractéristique $p$). L'argument de la seconde ligne de
la note est lu $G$ ou $G^{*}$ sans certitude. Cette formule est dans
Illusie, \emph{Complexe cotangent et déformations}~II (1972), chapitre
VII, cité de mémoire.}.

\emph{Lien avec la théorie de Dieudonné} (page 93). Sur un corps parfait,
si $M = D^{*}(G)$, la page conjecture — tout est au conditionnel — que
la suite exacte « du serpent » de $p = V F$ sur $M$,
\[
0 \to ({}_F M)^{(p)} \to ({}_p M)^{(p)} \xrightarrow{F} {}_V M \to M_F
\xrightarrow{V} (M_p)^{(p)} \to (M_V)^{(p)} \to 0,
\]
où ${}_F M = \mathrm{Ker}\, F$ et $M_F = \mathrm{Coker}\, F$, s'identifie
terme à terme à
\[
0 \to n_G \to D^{*}(G(1))^{\vee} \to t_{G^{*}} \to \omega_G \to
D^{*}(G(1)) \to \nu_{G^{*}} \to 0 ,
\]
et plus précisément que $\ell^{G} \simeq \mathcal{D}/F\mathcal{D}
\otimes^{L}_{\mathcal{D}} M$, avec $\mathcal{D}/F\mathcal{D} \simeq
k_{\sigma}[V]$. Quatre des identifications se contrôlent sur les groupes
d'ordre $p$ : $\omega_G \simeq M/FM$, $t_{G^{*}} \simeq \mathrm{Ker}\, V$,
$n_G \simeq \mathrm{Ker}\, F$ et $\nu_{G^{*}} \simeq \mathrm{Coker}\, V$,
à une torsion par $\sigma$ près\footnote{Contrôle de nous, sur $\mathbb{Z}/p$, $\mu_p$, $\alpha_p$.
La page note $G_p$ ce que nous écrivons $G(1)$ ; les deux premiers termes
de la ligne du haut sont lus avec doute.}. Elle propose enfin un
isomorphisme de triangles distingués entre
$\ell^{G}[-1] \to \Delta^{*}(G) \to \check{\ell}^{G^{*}}$ et le triangle
$\mathcal{D}/F\mathcal{D} \otimes^{L} M \to \mathcal{D}/p\mathcal{D}
\otimes^{L} M \to \mathcal{D}/V\mathcal{D} \otimes^{L} M$ déduit de $p =
VF$, pour un objet $\Delta^{*}(G)$ d'amplitude parfaite $[0, 1]$ qui
« devrait exister sur \emph{toute} base »\footnote{Le sommet du second
triangle est lu $\mathcal{D}/V\mathcal{D}$ sans certitude ; la flèche de
degré du premier porte « ? ».}. Cet objet hypothétique, avec $\omega_G$ en
sous-objet et $t_{G^{*}}$ en quotient, est ce que deviendra la
\emph{filtration de Hodge} du cristal de Dieudonné évalué sur $S$, que
la page 114 écrit $0 \to \omega_G \to \mathbb{D}(G)_S \to t_{G^{*}} \to
0$\footnote{Le rapprochement est de nous.}.

\pagerange{94}{94}
\subsection*{(C) Déformations de groupes commutatifs plats (page 94)}

La page attribue cette partie à Deligne et Illusie et renvoie en marge à
« ma lettre à Illusie ». Soit $S_0 = V(\mathcal{J}) \subset S$,
$\mathcal{J}^{2} = 0$. Les $\mathrm{Ext}$ qui suivent sont ceux de
faisceaux abéliens sur $S_0$, à valeurs dans des complexes de
modules\footnote{La page les indexe par $\mathcal{O}_{S_0}$, sur un
premier indice raturé ; $G_0$ n'est pas un module, et il s'agit d'$\mathrm{Ext}$
de faisceaux abéliens.}.

\begin{itemize}
  \item[(0)] Déformations des torseurs sous un groupe plat de présentation
  finie : gouvernées par $\mathrm{Ext}^{0}$ et $\mathrm{Ext}^{1}$ de
  $\ell^{G}$ dans $\mathcal{J}$.
  \item[(1)] Prolongement d'un homomorphisme $u_0 \colon G_0 \to H_0$ de
  groupes plats en $u \colon G \to H$, $G$ et $H$ relevant $G_0$ et $H_0$
  : obstruction dans $\mathrm{Ext}^{1}(G_0, \check{\ell}^{H_0}
  \otimes^{L} \mathcal{J})$, prolongements formant un torseur sous
  $\mathrm{Ext}^{0}(G_0, \check{\ell}^{H_0} \otimes^{L} \mathcal{J})$.
  \item[(2)] Déformation d'un groupe plat $G_0$ : obstruction dans
  $\mathrm{Ext}^{2}(G_0, \check{\ell}^{G_0} \otimes^{L} \mathcal{J})$,
  classes de déformations formant un torseur sous $\mathrm{Ext}^{1}$,
  automorphismes d'une déformation donnés par $\mathrm{Ext}^{0} =
  \mathrm{Hom}(G_0, \check{\ell}^{G_0} \otimes^{L} \mathcal{J})$\footnote{La
  page laisse en blanc les arguments de b) et c) et ajoute à côté,
  reliée par une accolade, « $(\underline{\mathrm{Hom}}(\underline{\omega}_{G_0},
  \mathcal{J}))$ », qui est $t_{G_0} \otimes \mathcal{J}$ lorsque
  $\omega_{G_0}$ est localement libre.}.
\end{itemize}
Une variante est annoncée pour les modules sur un anneau
$\mathcal{R}$\footnote{Lettre lue $\mathcal{R}$.}.

\pagerange{94}{95}
\subsection*{(D) Groupes de BT tronqués (pages 94 et 95)}

Renvoi : « Lettre Illusie 5.3 ».

\emph{Proposition.} Soit $G(n)$ un groupe de BT tronqué d'échelon $n$ sur
$S$, avec $p^{N} 1_S = 0$ et $n \geq N$. Alors $\omega_{G(n)}$ et
$n_{G(n)}$ sont localement libres de même rang. Pour $n \geq n' \geq N$,
l'inclusion $G(n') \hookrightarrow G(n)$ induit un isomorphisme
$\omega_{G(n)} \xrightarrow{\sim} \omega_{G(n')}$, et, si $n' \leq n - N$,
l'homomorphisme nul $n_{G(n)} \to n_{G(n')}$ ; la projection
$p^{n-n'} \colon G(n) \to G(n')$ induit un isomorphisme $n_{G(n')}
\xrightarrow{\sim} n_{G(n)}$ et l'homomorphisme nul $\omega_{G(n')} \to
\omega_{G(n)}$. Si $n \geq 2N$, $\omega_{G(n)} \simeq
n_{G(n)}$\footnote{La page écrit $H_i(L^{G(n)})$, que nous lisons comme
$H_i(\ell^{G(n)})$, seule lecture conforme à la notation de (B). Les
flèches vont dans le sens contravariant, comme sur la page. Nous ne
vérifions pas les bornes ; la démonstration n'est pas dans le dossier.}.

\emph{Théorème} (page 95). Soit $G$ un groupe de BT tronqué d'échelon $n$
et $M$ un module quasi-cohérent annulé par $p^{n}$. Alors, avec
$\Lambda_n = \mathbb{Z}/p^{n}\mathbb{Z}$,
\[
\underline{\mathrm{Ext}}^{1}_{\Lambda_n}(G, M) = 0, \qquad
\underline{\mathrm{Ext}}^{2}_{\Lambda_n}(G, M) \xrightarrow{\ \sim\ }
\underline{\mathrm{Ext}}^{2}_{\mathbb{Z}}(G, M).
\]
La page ajoute que $\underline{\mathrm{Ext}}^{2}_{\mathbb{Z}}(G, M) = 0$
pour tout $G$ fini localement libre si $2$ est inversible, « ce qui
simplifie la théorie des déformations pour les BT en car.\ $p \neq 2$
! »\footnote{« Inversible » est lu avec doute. Nous ne vérifions pas cette
annulation, et la laissons comme sienne. La parenthèse ouvrante manque sur
la page.}. Avec (C)(2), c'est la voie annoncée page 90 : les obstructions
et les classes de déformations d'un tronqué se lisent dans des
$\mathrm{Ext}$ que ce théorème contrôle.

\pagerange{96}{107}
\section*{VII. Les objets de Tate (pages 96 à 107)}

La page 96 est une chemise : « Sorites sur Groupes $p$-divisibles et
groupes de B.T. ». Les pages qui suivent isolent une structure abstraite
— un objet muni d'un endomorphisme nilpotent qui se comporte comme un
module libre sur $\mathbb{Z}/p^{n}$ — puis l'appliquent aux schémas en
groupes.

\pagerange{97}{101}
\subsection*{Dans une catégorie abélienne (pages 97, 99 et 101)}

Soit $\mathcal{C}$ une catégorie abélienne, $A$ un objet, $f$ un
endomorphisme de $A$ et $n \geq 1$ avec $f^{n} = 0$. On pose
$\mathrm{gr}^{i}_f(A) = f^{i}A / f^{i+1}A$ ; $f^{i}$ induit des
épimorphismes
\[
\theta^{i} \colon \mathrm{gr}^{0}_f(A) \twoheadrightarrow
\mathrm{gr}^{i}_f(A), \qquad \text{d'où}\quad
\theta \colon \mathrm{gr}^{0}_f(A)[t] \twoheadrightarrow
\mathrm{gr}^{\bullet}_f(A).
\]

\emph{Proposition.} Les conditions suivantes sont équivalentes :
\begin{itemize}
  \item[(i)] $\theta^{i}$ est un isomorphisme pour $i \leq n-1$ ;
  \item[(i bis)] $\theta^{i}$ est un monomorphisme pour $i \leq n-1$ ;
  \item[(i ter)] $\theta^{n-1}$ est un monomorphisme ;
  \item[(ii)] $(f^{i})^{-1}(f^{i+1}A) = fA$ pour $0 \leq i \leq n-1$ ;
  \item[(ii bis)] $\mathrm{Ker}\, f^{n-1} = fA$ ;
  \item[(iii)] $\mathrm{Ker}\, f^{j} = \mathrm{Im}\, f^{n-j}$ pour
  $0 \leq j \leq n$ ;
  \item[(iv)] il existe $j$, $1 \leq j \leq n-1$, tel que
  $\mathrm{Ker}\, f^{j} = \mathrm{Im}\, f^{n-j}$.
\end{itemize}
Un couple $(A, f)$ qui les satisfait est un \emph{objet de Tate de degré
$n$}. Exemple : un $\mathbb{Z}/p^{n}$-module libre, avec $f = p$.

\emph{Démonstration.} (i) $\Leftrightarrow$ (i bis) parce que les
$\theta^{i}$ sont épimorphiques. Le noyau de $\theta^{i}$ est
$(f^{i})^{-1}(f^{i+1}A)/fA$, d'où (i bis) $\Leftrightarrow$ (ii), et pour
$i = n-1$, (i ter) $\Leftrightarrow$ (ii bis). (i bis) $\Rightarrow$ (i
ter) est trivial, et (i ter) $\Rightarrow$ (i bis) parce que
$\theta^{n-1}$ est le composé de $\theta^{i}$ et de l'épimorphisme
$\mathrm{gr}^{i} \to \mathrm{gr}^{n-1}$ induit par
$f^{n-1-i}$\footnote{La page écrit $\theta^{n-1} =
\theta^{n-1-i}\theta^{i}$, abus de notation pour ce composé. Le
diagramme d'implications de la page est surchargé ; on en donne ici une
chaîne complète, dont la démonstration (i bis) $\Leftrightarrow$ (ii) est
de nous.}.

(ii bis) $\Rightarrow$ (iii). Toujours $\mathrm{Im}\, f^{n-j} \subset
\mathrm{Ker}\, f^{j}$. Pour l'inverse, soit $1 \leq j \leq n-1$ : on a
$\mathrm{Ker}\, f^{j} \subset \mathrm{Ker}\, f^{n-1} = fA$, donc tout
élément de $\mathrm{Ker}\, f^{j}$ est $fy$ avec $f^{j+1}y = 0$, soit
$\mathrm{Ker}\, f^{j} = f(\mathrm{Ker}\, f^{j+1})$, et de proche en
proche
\[
\mathrm{Ker}\, f^{j} = f\,\mathrm{Ker}\, f^{j+1} = \dots =
f^{n-1-j}\,\mathrm{Ker}\, f^{n-1} = f^{n-j}A
\]
\footnote{On raisonne « sur les éléments », ce que justifie le théorème
de plongement de Freyd–Mitchell, ou une réécriture en sous-objets. La page
présente cet argument comme preuve de (ii) $\Rightarrow$ (iii) ; il n'use
que de (ii bis).}. (iii) $\Rightarrow$ (ii bis) : faire $j = n-1$.
(iii) $\Rightarrow$ (iv) est trivial ($n \geq 2$).

(iv) $\Rightarrow$ (ii bis). Notons $S_j$ l'\emph{inclusion}
$\mathrm{Ker}\, f^{j} \subset \mathrm{Im}\, f^{n-j}$. Si $2 \leq j \leq
n-1$, $S_j$ entraîne $S_{j-1}$ : $\mathrm{Ker}\, f^{j-1} \subset
\mathrm{Ker}\, f^{j} \subset \mathrm{Im}\, f^{n-j} \subset \mathrm{Im}\, f$,
donc $\mathrm{Ker}\, f^{j-1} = f(\mathrm{Ker}\, f^{j}) \subset
f(\mathrm{Im}\, f^{n-j}) = \mathrm{Im}\, f^{n-j+1}$. De proche en proche,
$S_1$ : $\mathrm{Ker}\, f \subset f^{n-1}A$, et comme l'inclusion inverse
est toujours vraie, $\mathrm{Ker}\, f = f^{n-1}A$. Alors $f$ induit un
isomorphisme $A/f^{n-1}A \xrightarrow{\sim} fA$, qui envoie la filtration
$(f^{i}A/f^{n-1}A)_i$ \emph{sur} la filtration $(f^{i+1}A)_i$ ; c'est donc
un isomorphisme sur les gradués, c'est-à-dire que $f \colon
\mathrm{gr}^{i} \to \mathrm{gr}^{i+1}$ est un isomorphisme pour $i \leq
n-2$, et les $\theta^{i}$, qui en sont les composés, sont des
isomorphismes : c'est (i)\footnote{On remarque — c'est de nous — que la
démonstration n'utilise que l'inclusion $S_j$ pour un seul $j$ : elle
suffit déjà. La page conclut par un mot illisible après « c'est ».}.

\emph{Rang.} Supposons donnée sur $\mathcal{C}$ une fonction $\mathrm{rg}$ à
valeurs entières \emph{positives}, additive sur les extensions et nulle
seulement sur l'objet nul\footnote{La positivité est ajoutée par nous ;
la page dit « à valeurs entières », lu avec doute, et l'argument
l'exige.}. Comme $\mathrm{rg}(A) = \sum_{i < n} \mathrm{rg}\,
\mathrm{gr}^{i} \leq n\, \mathrm{rg}(A/fA)$, avec égalité si et seulement
si chaque $\theta^{i}$ est un isomorphisme,
\[
(A, f) \text{ est de Tate de degré } n \iff \mathrm{rg}(A) = n\,
\mathrm{rg}(A/fA).
\]

\emph{Morphismes} (pages 99 et 101). Soit $(B, g)$ un second objet de
Tate, de même degré $n$, et $u \colon A \to B$ avec $u f = g u$ et
$\mathrm{rg}(A/fA) = \mathrm{rg}(B/gB)$\footnote{La page écrit « la
condition analogue pour $(g, m)$ » ; l'argument demande $m = n$, que nous
supposons. Le passage est très raturé, et la liste de conditions qui
termine la page 99 n'est lisible que par fragments ; la page 101, dont la
transcription note qu'elle commence « en cours de démonstration », en est
la suite : la page 100 est le recto étranger de la feuille. Ce
raccord est de nous.}. Alors les conditions suivantes sont
équivalentes : $u$ est un isomorphisme ; $u$ est un monomorphisme ; $u$
est un épimorphisme ; $\mathrm{gr}^{0}(u)$ est un monomorphisme ;
$\mathrm{gr}^{0}(u)$ est un épimorphisme. En effet, si $u$ est mono,
$\mathrm{gr}^{0}(u)$ l'est, car $\theta^{n-1}$ l'identifie à
$f^{n-1}A \to g^{n-1}B$ ; si $u$ est épi, $\mathrm{gr}^{0}(u) \colon A/fA
\to B/gB$ l'est ; un mono ou un épi entre objets de même rang est un
isomorphisme ; si $\mathrm{gr}^{0}(u)$ est un isomorphisme, chaque
$\mathrm{gr}^{i}(u)$ l'est, puisque les $\theta^{i}$ l'identifient à
$\mathrm{gr}^{0}(u)$ ; et un morphisme de filtrations finies qui est un
isomorphisme sur les gradués est un isomorphisme.

\emph{Corollaire.} Pour $0 \leq i \leq n$, $f^{i}A$ est le seul
sous-objet de Tate de $A$ stable par $f$, de degré $n - i$ et de rang
$\rho = \mathrm{rg}(A/fA)$ : un tel $B$ est contenu dans $\mathrm{Ker}\,
f^{n-i} = f^{i}A$, qui est de Tate de degré $n - i$ et de rang $\rho$,
et l'inclusion, mono entre objets de Tate de même degré et de même rang,
est une égalité. \emph{Corollaire dual} : $A/f^{n-i}A$ est le seul quotient
de Tate de $A$ de degré $n - i$ et de rang $\rho$.

\pagerange{101}{103}
\subsection*{Parmi les schémas en groupes (pages 101 et 103)}

Soit $G$ un schéma en groupes commutatif, plat et de présentation finie
sur $S$\footnote{« Plat » est lu avec doute, et « loc.\ » est barré
devant « de présentation finie » ; la platitude de $G$ est nécessaire à
(ii) $\Rightarrow$ (i).}, $f$ un endomorphisme de $G$ et $n$ un entier
avec $f^{n} = 0$. Fixons $1 \leq j \leq n-1$. Les conditions suivantes
sont équivalentes :
\begin{itemize}
  \item[(i)] $\mathrm{Ker}\, f^{j}$ est plat sur $S$, et $f^{n-j} \colon G
  \to \mathrm{Ker}\, f^{j}$ est fidèlement plat ;
  \item[(ii)] pour tout $s \in S$, $(G_s, f_s)$ est un objet de Tate de
  degré $n$ dans la catégorie abélienne des groupes algébriques
  commutatifs sur $k(s)$.
\end{itemize}
\emph{Démonstration.} (i) $\Rightarrow$ (ii) : les noyaux commutent au
changement de base, et (i) donne $\mathrm{Ker}\, f_s^{j} = \mathrm{Im}\,
f_s^{n-j}$ ; on conclut par (iv) ci-dessus. (ii) $\Rightarrow$ (i) : soit
$H = \mathrm{Ker}\, f^{j}$, de présentation finie, et $g \colon G \to H$
induit par $f^{n-j}$ (il existe puisque $f^{j}f^{n-j} = 0$). Par (ii),
$g_s$ est fidèlement plat pour tout $s$ ; comme $G$ est plat de
présentation finie et $H$ de présentation finie, le critère de platitude
par fibres (EGA~IV, 11.3.10) montre que $g$ est fidèlement plat et que
$H$ est plat.

\emph{Corollaires.} La condition ne dépend pas de $j$ ; tous les
$\mathrm{Ker}\, f^{j}$ sont alors plats, et pour $j \geq j'$ les quotients
$\mathrm{Ker}\, f^{j}/\mathrm{Ker}\, f^{j'}$ sont représentables par des
schémas en groupes de présentation finie\footnote{La ligne qui suit, «
$f$ dans $G$ faisant … de la catégorie de Tate de hauteur $n$ (loc.\
cit.) », est surchargée et en partie illisible.}. Si $G$ est fini sur
$S$, le rang de Tate se lit en un point, et il y a un seul sous-objet de
Tate de degré donné $i \leq n$ et de rang $\rho$ donné ; « kif-kif pour
quotients ».

Pour $G$ fini localement libre, $f = p\,\mathrm{id}_G$ et $n \geq 2$, ce
sont exactement les groupes de BT tronqués d'échelon $n$ : un groupe fini
localement libre annulé par $p^{n}$ avec $\mathrm{Ker}\, p^{j} = \mathrm{Im}\, p^{n-j}$\footnote{Le
rapprochement est de nous ; pour $n = 1$ la définition des groupes
tronqués demande en plus $\mathrm{Ker}\, F = \mathrm{Im}\, V$, ce que les
objets de Tate de degré $1$ ne voient pas.}. Les pages 97 à 103 donnent
donc une théorie des groupes de BT tronqués qui ne sait d'eux que leur
endomorphisme $p$.

\pagerange{105}{107}
\subsection*{Un sous-groupe d'un objet de degré 2 (pages 105 à 107)}

Le cadre, que la page ne dit pas, se reconstitue : $S$ est un trait de
point générique $\eta$, $G_2$ un groupe fini plat sur $S$ qui est de Tate
de degré $2$ pour $p$, $G_1 = \mathrm{Ker}\, p$, de sorte que $p$ induit
$G_2/G_1 \xrightarrow{\sim} G_1$ ; $G'_2 \subset G_2$ est un sous-groupe
fini plat, adhérence schématique de sa fibre générique, et $G''_2 =
G_2/G'_2$\footnote{Reconstitution de nous, d'après les formules
$G'_2 = \overline{G'_{2\eta}}$ et l'emploi de $p$.}. On pose $G'_1 =
\overline{G'_{2\eta} \cap G_{1\eta}}$. Alors
\[
G'_1 \subset \overline{G'_{2\eta}} \cap \overline{G_{1\eta}} = G'_2 \cap G_1
= \mathrm{Ker}(p|_{G'_2}),
\]
avec égalité si et seulement si $G'_2 \cap G_1$ est plat sur $S$. La
multiplication par $p$ envoie $G'_2$ dans $G'_1$ (fibre par fibre, puis
par platitude) et tue $G'_1$, d'où $G'_2/G'_1 \to G'_1$, de noyau
$(G'_2 \cap G_1)/G'_1$. C'est aussi le noyau de la flèche $\alpha'' \colon
G''_1 = G_1/G'_1 \to G''_2$ induite par les inclusions $G'_1 \subset G'_2$,
$G_1 \subset G_2$, et la page l'inscrit dans la suite exacte
\[
0 \to \mathrm{Ker}\, \alpha'' \to G'_2/G'_1 \to G_2/G_1 \to
G''_2/\mathrm{Im}\, G''_1 \to 0
\]
\footnote{La page dessine un morphisme de suites exactes courtes entre
$0 \to G'_2 \to G_2 \to G''_2 \to 0$ et $0 \to G'_1 \to G_1 \to G''_1 \to
0$, flèches verticales $\alpha'$, $\alpha$, $\alpha''$ tracées de la
première ligne vers la seconde. Le noyau qu'elle calcule pour $\alpha''$,
$(G'_2 \cap G_1)/G'_1$, et le conoyau $G''_2/\mathrm{Im}\, G''_1$ sont
ceux de la flèche induite par les inclusions, qui va de $G''_1$ vers
$G''_2$ ; nous l'orientons ainsi (lecture de nous), et ne reproduisons
pas le diagramme. En tête de page, « $G' \subset G \subset G''$ ».}.

Page 107, la liste de conditions équivalentes : a) $G'_1 = G'_2 \cap G_1$ ;
b) $G'_2 \cap G_1$ est plat ; c) $G'_2$ est de Tate de degré $2$ ; d)
$G'_2/G'_1 \to G'_1$ est un monomorphisme ; e) c'en est un épimorphisme ;
f) $p \colon G'_2 \to G'_1$ est un épimorphisme. Les équivalences a)
$\Leftrightarrow$ b) $\Leftrightarrow$ d) valent toujours. Celles qui
font intervenir c), e), f) demandent que $G'_{2\eta}$ soit elle-même de
Tate de degré $2$ — ce qui donne à $G'_2/G'_1$ et $G'_1$ le même rang —,
et sous cette hypothèse les six conditions sont
équivalentes\footnote{L'hypothèse et le contre-exemple sont de nous :
$G_2 = \mathbb{Z}/p^{2}$, $G'_2 = p\mathbb{Z}/p^{2} \simeq \mathbb{Z}/p$
satisfont a) et d), pas c) ni e). La page a écrit « mono » pour d) et,
sous des hachures, « épim ».}. La page commence une liste duale pour le
quotient — f*) $G''_1 \to G''_2$ mono, c*) $G''_1$ de Tate de degré $2$ —
et la laisse inachevée.

Un encadré, en dessous, rassemble des essais d'exemples : un groupe dont
la fibre générique est $(\mathbb{Z}/p^{2} \times \mathbb{Z}/p)_\eta$,
$\mathbb{Z}/p^{2} \to \alpha_p$, et un groupe $C$ muni de $u \colon C \to
C$ relevant $F$, avec les relations formelles $V = uF$, $F = u^{-1}V$,
$V^{2} = up$, $F^{2} = u^{-1}p$ tirées de $VF = p$\footnote{$V^{2} = up$
suppose que $u$ commute à $F$ ; la page ne le dit pas. Le but de ces
essais n'est pas écrit.}.

La page 106, brouillon dans plusieurs sens, reprend les calculs de
$\mathrm{Ker}\, f^{j} = \mathrm{Im}\, f^{n-j}$ (« $\mathrm{Ker}\, f =
\mathrm{Im}\, f^{n-1} \Rightarrow \mathrm{Ker}\, f^{j} = \mathrm{Im}\,
f^{n-j}$ », qui est l'étape $S_1$ ci-dessus). Elle s'ouvre sur quatre
lignes en anglais, d'une main que le feuillet ne permet pas d'attribuer :
« ? Proof of 1st conj. : given $0 \to \Phi' \to \Phi(G) \to \Phi'' \to
0$ (Galois-devissage), then can define $G'$, $G''$ so $\Phi' =
\Phi(G')$, $\Phi'' = \Phi(G'')$ ». Ni la conjecture ni $\Phi$ ne sont
définis.

\pagerange{109}{109}
\section*{VIII. Épimorphismes formellement lisses (page 109)}

Pour une suite exacte $0 \to N \to G \xrightarrow{u} H \to 0$, la page
note d'abord qu'elle n'entraîne pas que $N$ soit formellement lisse. Elle
cherche ensuite des conditions équivalentes pour qu'un homomorphisme $u$
soit formellement lisse et épimorphique : (iii) $u$ l'est ; (iii bis)
$u_s$ l'est sur la fibre spéciale ; (iv) $G_k \to H_k$ est un épimorphisme
(formellement lisse) et $\Phi(G) \to \Phi(H)$ est surjectif ; (i) et
(ii), sur les points à valeurs dans un objet que la page ne laisse pas
lire, sont illisibles en partie, et (iv bis) est vide. Deux remarques :
$\Phi$ est « conservatif » — si $\Phi(u)$ est un isomorphisme, $u$ en est
un — et, si $u$ est formellement fidèlement plat sans être formellement
lisse, $\Phi(u)$ n'est pas surjectif\footnote{Le foncteur $\Phi$ n'est
défini nulle part dans le dossier ; une ligne barrée,
« $G(\hat{S}) \to H(\hat{S})$ », suggère des points à valeurs dans une
complétion, et la « Galois-devissage » de la page 106 un module
galoisien, sans que rien ne tranche. Nous ne reconstituons pas les
conditions illisibles.}.

\pagerange{112}{114}
\section*{IX. « Conférence à Nice » (pages 112 à 114)}

Trois pages au propre, titrées de sa main, encadré : le canevas d'un
exposé. Il est vraisemblable qu'il s'agisse de sa conférence au Congrès
international des mathématiciens de Nice, en septembre 1970, publiée
sous le titre « Groupes de Barsotti–Tate et cristaux »\footnote{Identification
de nous, que nous n'avons pas comparée au texte imprimé (\emph{Actes du
Congrès international des mathématiciens}, 1970, t.~1), cité de
mémoire. Elle s'accorde avec la datation de l'inventaire.}. Le plan va du
groupe de BT au foncteur de Dieudonné cristallin, en quatre parties.

\pagerange{112}{112}
\subsection*{1) et 2) Groupes de Barsotti–Tate, groupe formel (page 112)}

Topologie fppf sur $S$, $p$ fixé. $G$ est un groupe de BT si a) $pG = G$,
b) $G$ est de $p$-torsion, c) $G(1)$ est fini localement libre. Alors
chaque $G(n)$ l'est, de rang $p^{nh}$ si $G(1)$ est de rang $p^{h}$, et
$G = \varinjlim G(n)$\footnote{La page note $d$ ce que nous notons $h$,
la hauteur ; ailleurs $d$ est la dimension (page 89).}. Exemple
important : pour un schéma abélien $A$, $G = A(p^{\infty})$,
$G(n) = A(p^{n})$.

\emph{Théorème de Serre–Tate.} Si $p$ est localement nilpotent sur $S$ et
$S_0 \hookrightarrow S$ défini par un idéal nilpotent, le foncteur
$A \mapsto (A_0, A(p^{\infty}), \varphi)$, où $\varphi$ identifie
$A(p^{\infty})_{S_0}$ à $A_0(p^{\infty})$, est une équivalence entre
schémas abéliens sur $S$ et triplets formés d'un schéma abélien $A_0$ sur
$S_0$, d'un groupe de BT sur $S$ et d'un tel
isomorphisme\footnote{Hypothèses ajoutées ; la page écrit, sous
$S_0 \hookrightarrow S$, quelques mots lus avec doute (« défini de »,
« com.\ $p$ »). Le théorème est de Serre et Tate (1964) ; une
démonstration publiée est celle de Katz (1981), d'après Drinfeld, citée
de mémoire.}. Vient ensuite le groupe formel $\widehat{G}$ associé à un
groupe de BT (le théorème~2.5 de la page 89).

\pagerange{112}{113}
\subsection*{3) Cristaux (pages 112 et 113)}

\emph{Puissances divisées.} Sur un idéal $I$, des opérations $x \mapsto
x^{[n]}$ ($n \geq 0$), qui miment $x^{n}/n!$ : $x^{[0]} = 1$, $x^{[1]} =
x$, $x^{[n]} \in I$ pour $n \geq 1$,
\[
(x+y)^{[n]} = \sum_{i+j=n} x^{[i]} y^{[j]}, \qquad
(\lambda x)^{[n]} = \lambda^{n} x^{[n]},
\]
et les autres identités usuelles\footnote{La page écrit $x^{(n)}$ et
s'arrête sur « … » : $x^{[m]}x^{[n]} = \binom{m+n}{m} x^{[m+n]}$ et
$(x^{[m]})^{[n]} = \frac{(mn)!}{(m!)^{n}\, n!}\, x^{[mn]}$.}. On note
$I^{[n]}$ l'idéal engendré par les monômes de poids $\geq n$ en les
$x^{[i]}$. Les puissances divisées sont \emph{nilpotentes} si $I^{[n]} = 0$
pour $n$ grand, et la page appelle \emph{quasi-nilpotentes} celles pour
lesquelles $(n-1)!\, I^{[n]} = 0$ pour $n$ grand. L'exponentielle $\exp(x)
= \sum_{n \geq 0} x^{[n]}$ a un sens dans le premier cas, le logarithme
$\log(1+x) = \sum_{n \geq 1} (-1)^{n-1} (n-1)!\, x^{[n]}$ dans le
second\footnote{$(n-1)!\, x^{[n]}$ est l'écriture divisée de $x^{n}/n$.
Le mot « quasi-nilpotent » n'est pas, en ce sens, la terminologie
d'aujourd'hui.}.

\emph{Site cristallin} (absolu, ou relatif) : ses objets sont des
épaississements à puissances divisées, qu'on exige compatibles aux
puissances divisées usuelles de $p$\footnote{La fin de la phrase est
coupée ; on lit en dessous « nilpot. » et « Berthelot ».}. En marge :
« début de théorie cohomologique (cf.\ Illusie) développée par
Berthelot »\footnote{La thèse de Berthelot, \emph{Cohomologie cristalline
des schémas de caractéristique $p > 0$}, paraît en 1974 (cité de
mémoire).}.

Les \emph{cristaux de modules} localement libres sont les modules
localement libres sur le site cristallin dont les morphismes de
transition sont des isomorphismes. Deux exemples :
\begin{itemize}
  \item si $S = \mathrm{Spec}\, A_0$, $A_0$ parfait de caractéristique
  $p$, les cristaux sur $A_0$ sont les modules localement libres sur
  $A = W(A_0)$ ;
  \item si $X_0$ est lisse sur un corps parfait $k$, $W = W(k)$, et
  $\mathcal{X}$ un schéma formel lisse sur $W$ relevant $X_0$, un cristal
  est un module localement libre $M$ sur $\mathcal{X}$ muni d'une
  connexion intégrable (à courbure nulle) relativement à $W$, et
  topologiquement quasi-nilpotente\footnote{Cette dernière condition,
  nécessaire, est ajoutée par nous ; la page ne la dit pas. Elle barre
  un ajout sur un relèvement $F$ du Frobenius, $\sigma_W$-linéaire.}.
\end{itemize}
« Ce sont des genres d'animaux rencontrés naturellement par Dwork,
Katz, … »

\pagerange{113}{114}
\subsection*{4) Le programme (pages 113 et 114)}

\begin{itemize}
  \item[I)] Pour $S$ un schéma où $p$ est localement nilpotent, on définit
  un « foncteur de Dieudonné »
  \[
  \mathrm{BT}(S)^{\circ} \longrightarrow \{\text{cristaux en modules
  localement libres sur } S\}, \qquad G \longmapsto \mathbb{D}(G),
  \]
  compatible au changement de base.
  \item[II)] Si $S$ est de caractéristique $p$, $M = \mathbb{D}(G)$ est
  muni de $F_M$, $V_M$ avec $F_M V_M = V_M F_M = p\,\mathrm{id}$ : c'est
  un \emph{cristal de Dieudonné}. « Il se pourrait que le foncteur
  $\mathrm{BT}(S)^{\circ} \to \{\text{cristaux de Dieudonné}\}$ soit
  pleinement fidèle, et j'ai une idée de ce que pourrait être l'image
  essentielle… »\footnote{La pleine fidélité n'est démontrée aujourd'hui
  que sous des hypothèses sur la base — de Jong (1995) pour les bases
  formellement lisses sur un corps parfait, et d'autres classes depuis ;
  références citées de mémoire. Le dossier ne dit pas quelle image
  essentielle il avait en vue.}.
  \item[III)] Pour $S = \mathrm{Spec}\, k$, $k$ parfait de
  caractéristique $p$, on retrouve le foncteur de Dieudonné classique de
  la page 81.
  \item[IV)] $\mathbb{D}(G)_S$ porte une filtration
  \[
  0 \to \omega_G \to \mathbb{D}(G)_S \to t_{G^{*}} \to 0
  \]
  \footnote{L'exposant du premier $\omega$ est surchargé ; nous lisons
  $\omega_G$, la reconstruction est de nous et s'accorde avec la page 93
  et avec $\check{\omega}_{G^{*}} = t_{G^{*}}$ au troisième terme.} ;
  c'est la filtration de Hodge. Soit $S_0 \hookrightarrow S$ défini par
  un idéal muni de puissances divisées (nilpotentes), $p$ localement
  nilpotent sur $S$. Alors prolonger $G_0$ en un groupe de BT sur $S$
  est « la même chose » que prolonger la filtration de
  $\mathbb{D}(G_0)_S$, c'est-à-dire relever $\omega_{G_0}$ en un
  sous-module localement facteur direct de $\mathbb{D}(G_0)_S$. La page
  le restreint aux $G_0$ dont les fibres sont toroïdales ou
  connexes\footnote{Lecture en partie douteuse (« composées »,
  « toroïdales »). C'est le théorème de Grothendieck–Messing, que Messing
  démontre en 1972 sans restriction sur les fibres, pour des puissances
  divisées nilpotentes.}.
  \item[V)] À isogénie près, sur un anneau $A$ noethérien complet pour la
  topologie $p$-adique, $A_0 = A/pA$, le foncteur $G \mapsto (G_0,
  \text{filtration sur } \mathbb{D}(G_0)_A)$ serait pleinement fidèle. Pour
  un anneau de valuation discrète complet à corps résiduel parfait $k$,
  un groupe de BT serait connu à isogénie près par le module de Dieudonné
  $M$ de sa fibre spéciale et une filtration sur $M \otimes_W A$, rendu
  rationnel\footnote{La page écrit « $(M \otimes_W A)_p$ », que nous
  lisons comme $M \otimes_W A[1/p]$. Cette description des groupes
  $p$-divisibles par des modules filtrés est l'objet d'une longue suite
  de travaux : Fontaine (1977), puis Breuil et Kisin pour l'énoncé
  général sur un anneau d'entiers $p$-adique — cités de mémoire.}.
\end{itemize}

Le dossier s'arrête sur ce programme. Ce qu'il a mis en place tient en
peu de mots : le groupe de BT et sa suite connexe-étale (pages 83 à 87),
son groupe formel (page 89), la théorie des déformations par le complexe
cotangent et l'annonce qu'un groupe de BT se déforme sans obstruction
(pages 90 à 95), une théorie abstraite des groupes tronqués (pages 97 à
107) — et, page 114, l'énoncé qui fait de la déformation d'un groupe une
affaire d'algèbre linéaire.

\end{document}
